\documentclass[11pt]{article}
\usepackage[margin=1in]{geometry}
\usepackage{amsmath,amssymb}
\usepackage{booktabs}
\usepackage{url}
\usepackage[hidelinks]{hyperref}

\newcommand{\Ro}{\mathrm{Ro}}
\newcommand{\Rb}{R_\beta}
\newcommand{\dpar}[2]{\frac{\partial #1}{\partial #2}}

\title{Fixed-$\Ro$ equal-area theory on the equatorial $\beta$ plane,\\
for the SS09 shallow-water model}
\author{pengcheng-ss09 suite, Tier 0 (A0 memo)}
\date{2026-07-16}

\begin{document}
\maketitle

\noindent
Beta-plane translations of the closed forms in the fixed-$\Ro$ manuscript
(\path|papers/fixed-ro-eq-area_ms-draft-v1.pdf|, ``ms'' below), derived
against this model's exact equations. Every result carries a basis tag;
``pending run'' tags get upgraded as the suite's runs confirm them.
Section~\ref{sec:checks} lists the falsifying checks and their status, and
the run record follows it.

\section{Setup and notation}

Steady state, hemispherically symmetric forcing ($y_0 = 0$), northern
hemisphere ($y > 0$). The model's steady equations (\texttt{SCIENCE.md}
section~2), with $T = \theta/1.6$ and the $v$ equation's factor of two
multiplying only the time tendency (so it drops out of every steady
balance):
\begin{align}
  v\left(\beta y - \dpar{u}{y}\right)
    &= \mathcal{H}(\theta_E - \theta)\, u \dpar{v}{y}
       + \epsilon_u u
       + v_d\, \mathcal{H}(u)\, \mathrm{sgn}(y) \dpar{u}{y},
    \label{eq:umom} \\
  \beta y\, u
    &= -\frac{g H}{T_0} \dpar{T}{y} + k_v \frac{\partial^2 v}{\partial y^2},
    \label{eq:vmom} \\
  \frac{\delta \Delta_z}{H} \dpar{v}{y}
    &= \frac{\theta_E - \theta}{\tau}.
    \label{eq:thermo}
\end{align}
Forcing (SS09 parabolic profile, the theory-matched choice):
$\theta_E = \theta_{00} - \Delta_y (y/y_1)^2$ for $|y| < y_1$. In
temperature units,
\begin{equation}
  T_E(y) = T_E(0) - b\, y^2,
  \qquad b \equiv \frac{\Delta_T}{y_1^2},
  \qquad \Delta_T \equiv \frac{\Delta_y}{1.6},
\end{equation}
with $T_0 = 300$~K. The small-angle limit of the ms is exact here: the
$\beta$ plane has no $\cos\varphi$ factors, so no truncation enters any of
the following.

\section{The one-line mapping}

Every ms result carries over with the substitutions
\begin{equation}
  \Ro_{\mathrm{th}} \;\to\; \Rb \equiv \frac{4 g H \Delta_T}{T_0 \beta^2 y_1^4},
  \qquad
  \Delta_h \theta_0 \;\to\; \Delta_T,
  \qquad
  a\varphi \;\to\; y,
  \qquad
  \varphi_{\Ro} \;\to\; Y_{\Ro}/y_1.
\end{equation}
\textbf{[derived]} $\Rb$ follows from the sphere definition
$\Ro_{\mathrm{th}} = g H \Delta_h / (\Omega a)^2$ with $\Omega = \beta a/2$
and $\Delta_h = \Delta_T a^2 / (T_0 y_1^2)$ (matching the forcing curvature
at the equator); the planetary radius cancels, as it must on a $\beta$
plane. At this repo's defaults ($g = 9.81$~m\,s$^{-2}$, $H = 16$~km,
$\Delta_y = 50$~K so $\Delta_T = 31.25$~K, $T_0 = 300$~K,
$\beta = 2\times 10^{-11}$~m$^{-1}$\,s$^{-1}$, $y_1 = 9439$~km):
\begin{equation}
  \boxed{\;\Rb = 0.0206.\;}
\end{equation}

\section{Wind and temperature}
\label{sec:wind}

\textbf{[derived: pending run confirmation]}
Uniform $\Ro = (\partial_y u)/(\beta y)$ integrates to
$u_{\Ro}(y) = \Ro\, \beta y^2/2$ (ascent at $y = 0$). Gradient balance
(neglecting $k_v$) gives $\partial_y T = -(T_0 / g H)\, \Ro\, \beta^2 y^3/2$,
so
\begin{equation}
  T_{\Ro}(y) = T_{\Ro}(0) - c\, y^4,
  \qquad c \equiv \frac{\Ro\, T_0 \beta^2}{8 g H}.
  \label{eq:TRo}
\end{equation}

\section{Equal-area edge and equatorial temperature}
\label{sec:edge}

\textbf{[derived: scalar anchors confirmed at the 1--8\% level by runs
1a--1b below; at profile level the inner half of the cell agrees at the
10\% level and the terminus misfit is ${\sim}35$--$40\%$ of the
depression, run-record addendum]}
With $D(y) = T_{\Ro}(y) - T_E(y)$, imposing $D(Y) = 0$ and
$\int_0^Y D\, \mathrm{d}y = 0$ (measure $\mathrm{d}y$; the Newtonian
relaxation has uniform $\tau$) gives $\tfrac{2}{3} b Y^2 = \tfrac{4}{5} c
Y^4$, i.e. $Y^2 = \tfrac{5}{6}\, b/c$, which in Section-2 variables is
\begin{align}
  Y_{\Ro} &= y_1 \left( \frac{5 \Rb}{3 \Ro} \right)^{1/2}
  &&\text{[ms Eq.~13 analog]},
  \label{eq:Yro} \\
  T_E(0) - T_{\Ro}(0) &= \frac{5}{18} \frac{\Rb\, \Delta_T}{\Ro}
  &&\text{[ms Eq.~15 analog]}.
  \label{eq:depression}
\end{align}
Defaults, $\Ro = 1$: $Y_{\Ro} = 1.75\times 10^6$~m and equatorial
depression $0.18$~K (in $T$; multiply by 1.6 for $\theta$ units). At
$\Ro = 0.15$ the depression is 1.2~K ($T$) / 1.9~K ($\theta$).

Consistency note: the analogous derivation for the sin$^2$ profile needs
the numerical equal-area solver (puffins \texttt{equal\_area\_bouss} with
$\Ro$-scaled \texttt{del\_h\_over\_ro}, verified in puffins PR~\#46). Near
the equator $\sin^2\!\big(\pi y/(2 y_1)\big) = \big(\pi y/(2 y_1)\big)^2 +
O(y^4)$, so sin$^2$ behaves like a parabola with curvature enhanced by
$(\pi/2)^2 = 2.47$; Section~\ref{sec:closure}'s postdiction check uses
exactly this factor.

\section{Meridional wind, heat flux, momentum flux}
\label{sec:v}

\textbf{[derived: $v_{\max}$ magnitude and location confirmed at the
1--4\% level by runs 1a--1b below; at profile level the inner flank agrees
at the 5\% level and the terminus flank misfit is ${\sim}45\%$ of
$v_{\max}$, run-record addendum]}
The steady thermodynamic balance \eqref{eq:thermo} integrates to the
model's $v$ profile directly, using $\theta_E - \theta = 1.6\,(T_E - T) =
-1.6\, D$:
\begin{equation}
  v(y)
  = -\frac{1.6 H}{\delta \Delta_z \tau} \int_0^y D\, \mathrm{d}y'
  = \frac{1.6 H}{\delta \Delta_z \tau}\,
    \frac{5}{18} \frac{\Rb \Delta_T}{\Ro}\,
    Y_{\Ro} \left( x - 2x^3 + x^5 \right),
  \qquad x \equiv \frac{y}{Y_{\Ro}}.
  \label{eq:v}
\end{equation}
The bracket is exactly the ms Eq.~(16) shape; it peaks at $x = 5^{-1/2} =
0.447$ with value $0.286$, and $v_{\max} \propto \Ro^{-3/2}$ (ms Eq.~20
analog). Defaults, $\Ro = 1$: $v_{\max} = 3.0\times 10^{-3}$~m\,s$^{-1}$ at
$y = 0.78\times 10^6$~m.

Column fluxes, using the two-layer reading (upper slab of depth $\delta$
with velocity $+v$ and potential temperature $\theta_{\mathrm{top}}$,
return slab of depth $\delta$ with $-v$ and $\theta_{\mathrm{bot}}$,
$\theta_{\mathrm{top}} - \theta_{\mathrm{bot}} = \Delta_z$; return-flow $u$
negligible):
\begin{align}
  \text{heat flux:} \quad
    &\int v\, \theta\, \mathrm{d}z = \delta\, \Delta_z\, v(y)
    &&\text{[ms Eq.~16]}, \\
  \text{momentum flux:} \quad
    &\int v\, u\, \mathrm{d}z = \delta\, v(y)\, u_{\Ro}(y)
    &&\text{[ms Eq.~17]}, \\
  \text{vertical velocity:} \quad
    &w = \delta \dpar{v}{y}, \qquad w_{\max} \propto \Ro^{-1}
    &&\text{[ms sec.~3f]}.
\end{align}
The ms's counterintuitive momentum-flux result carries over: expanding
$\delta v u_{\Ro}$, the leading term in $y$ is $\Ro$-independent (the
$\Ro^{-1}$ strengthening of $v$ cancels the $\Ro$ weakening of $u$ in the
product) and $\partial(\text{flux})/\partial \Ro < 0$ strictly inside the
cell, exactly as verified for the sphere forms in puffins
\texttt{test\_eq\_area.py}.

\section{Diagnosed quantities: conventions for the suite}
\label{sec:diag}

\begin{itemize}
  \item \textbf{Full-profile scorecard (the confirmation standard for
    every suite run):} the simulated $u$, $T$, and $v$ profiles are
    compared against the zero-free-parameter theory profiles at the
    diagnosed (regression) cell-mean $\Ro$, reporting $\max|\Delta|$ and
    its location both over the cell $[0, Y_{\Ro}]$ and over its inner half,
    each scaled by the field's theory amplitude (cell-edge $u$, the
    equatorial depression, and $v_{\max}$ respectively). The temperature
    comparison is visualized as the equal-area lobes $D(y) = T - T_E$,
    which removes the ${\sim}200$~K shared background. Implemented in
    \path|scripts/fixed_ro_scorecard.py| (one four-panel figure per run;
    fields are means over the final 30 days). Scalar anchors (depression,
    $v_{\max}$, edge) are by-products of the same tool, not the
    confirmation standard: runs 1a--1b passed the scalars at the 1--8\%
    level while the profile misfit at the terminus is ${\sim}35$--$45\%$ of
    scale (see the profile addendum in the run record).
  \item \textbf{Local $\Ro(y)$:} $(\partial_y u)/(\beta y)$, the model's
    existing diagnostic (NaN within $|y| < \Delta y_{\mathrm{grid}}$ of the
    equator).
  \item \textbf{Cell-mean $\Ro$, two definitions, both reported}
    (paralleling Hill et al.\ 2025): (i) the area mean of local $\Ro$ over
    $[y_{\Ro\text{-max}},\, Y_{\mathrm{edge}}]$; (ii) the best-fit $\Ro$
    regressing $u(y)$ against $\beta y^2/2$ over the same interval. The
    near-ascent region is excluded because $\Ro \to 0$ at the equator
    whenever $v_d > 0$ (Zhang et al.\ 2025, $u \sim y^3$).
  \item \textbf{Cell edge, three definitions, all reported:} the
    $v$-threshold edge ($|v|$ below 10\% of its extremum, existing
    diagnostic; see run 1a for its failure mode at small $v_d$), the jet
    latitude (existing spline diagnostic), and the $\theta$-merge latitude
    (where $\theta$ rejoins $\theta_E$ within a tolerance, searched
    poleward of the \emph{jet}; see run 1a). The theory says they coincide;
    their splitting is itself a result.
  \item \textbf{Momentum-budget shares:} each steady term of
    \eqref{eq:umom} evaluated from output; the drag share $\epsilon_u u$
    and the vertical-advection share quantify the accuracy of the
    EMFD--Coriolis balance underlying Zhang et al.\ (2025) Eq.~24,
    $\Ro = v/(v + v_d)$.
\end{itemize}

\section{Emergent cell-mean $\Ro$ closure}
\label{sec:closure}

\textbf{[derived: pending run confirmation]}
Combining Zhang's Eq.~24 at the cell-mean level with
Section~\ref{sec:v}'s $v$ magnitude (the cell mean of the bracket over
$[0,1]$ is $1/6$):
\begin{equation}
  \bar{\Ro} = \frac{\bar v}{\bar v + v_d},
  \qquad
  \bar v(\bar{\Ro})
  = \frac{1.6 H}{\delta \Delta_z \tau}\,
    \frac{5}{18} \frac{\Rb \Delta_T}{\bar{\Ro}}\,
    \frac{Y_{\Ro}(\bar{\Ro})}{6},
  \label{eq:closure}
\end{equation}
a transcendental equation for $\bar{\Ro}(v_d)$ with no simulation input.
This is the suite's analysis~G; it extends Zhang et al.\ (2025) section~4
(their max-$\Ro$ scalings $v_d^{-2/5}$, $v_d^{-1/4}$) to the cell-mean
variable the ms uses.

Postdiction check \textbf{[computed]}: Section~\ref{sec:v}'s $v$ with the
sin$^2$ profile's larger near-equator curvature ($(\pi/2)^2 = 2.47$ in
$b$, entering $v$ through $b^2 Y/c$ as a factor $\sim 9.6$) predicts
$v_{\max} \sim 0.03$~m\,s$^{-1}$ for the published defaults, matching
Zhang et al.\ (2025) Figs.~2--3 magnitudes within $\sim$1.5$\times$.

\section{Falsifying checks and status}
\label{sec:checks}

Every run in the checks below also gets the Section~\ref{sec:diag}
full-profile scorecard on $u$, $T$, and $v$; the tier-specific items are
the additional, targeted comparisons.

\begin{enumerate}
  \item \textbf{Full $u$/$T$/$v$ profiles, edge, and equatorial depression
    vs theory at diagnosed $\Ro$ (the most-falsifying single check):}
    SS09-parabolic forcing, $v_d = 0$. Framing correction from the first run: with the default
    vertical momentum advection on, $v_d = 0$ is \emph{not} the $\Ro = 1$
    limit (Zhang et al.\ 2025: $u \sim \beta y^2/3$ near the ascent, and
    drag lowers it further); the theory must be evaluated at the run's
    diagnosed cell-mean $\Ro$. The $\Ro \to 1$ anchor requires
    \texttt{-{}-no-vert-advec-u} \emph{and} reduced drag (run 1b). Status:
    runs 1a--1b below, scalar level and profile level (addendum).
  \item \textbf{Edge collapse:} $Y_{\mathrm{edge}} (\Ro/\Rb)^{1/2}/y_1 =
    (5/3)^{1/2}$ flat across the $v_d$, $\Delta_y$, $\beta$, and radius
    ladders. Status: pending Tiers 1--2.
  \item \textbf{Radius-ladder null:} dimensional $Y_{\mathrm{edge}}$
    invariant under $(\beta \to \beta/k,\; y_1 \to k y_1,\; L \to k L)$.
    Status: pending Tier 2.
  \item \textbf{$v$-profile shape and $\Ro^{-3/2}$ magnitude scaling}
    (Section~\ref{sec:v}). Status: shape checked at profile level in runs
    1a--1b (addendum); the $\Ro^{-3/2}$ scaling pending Tier 1.
  \item \textbf{Momentum flux increases as $\Ro$ decreases}
    (Section~\ref{sec:v}). Status: pending Tier 1.
  \item \textbf{Emergent $\bar{\Ro}(v_d)$} (Section~\ref{sec:closure}).
    Status: pending Tier 1 + analysis G.
\end{enumerate}

\section*{Run record}

Each confirming or refuting run gets an entry: configuration, output path,
measured vs predicted.

\subsection*{2026-07-16, check 1a: $v_d = 0$, default physics (vertical advection on)}

Config: SS09 parabolic forcing, $v_d = 0$, $n_y = 801$, $\Delta t = 30$~s,
numba backend, steady-state stop (window 30~d, threshold $10^{-4}$),
converged day 959. Output:
\path|model_output/fixed_ro_suite/tier0_amc_edge/amc_ss09_ny801_dt30_vd0.nc|.
Steadiness caveat \textbf{[computed]}: the jet was still drifting
$+0.4\%$ per 60~d at stop (drag-timescale tail), so numbers carry roughly
a 2\% allowance.

Diagnosed: cell-mean $\Ro$ 0.469 (area mean) / 0.470 ($u$ regression);
local $\Ro(y)$ is \emph{not} uniform: ${\sim}0.68$ at the ascent (the
vertical-advection $2/3$ of Zhang et al.\ 2025, drag-reduced), declining
through ${\sim}0.6$ mid-cell to ${\sim}0.2$ approaching the jet, i.e.\ the
ms's linear-$\Ro$ motivation in miniature.

Theory at diagnosed $\Ro = 0.470$ vs run \textbf{[computed]}:

\begin{center}
\begin{tabular}{lccc}
\toprule
quantity & predicted & measured & rel.\ gap \\
\midrule
equatorial depression $T_E(0) - T(0)$ & 0.381~K & 0.379~K & 0.5\% \\
$v_{\max}$ & $9.28\times 10^{-3}$~m/s & $9.39\times 10^{-3}$~m/s & 1.2\% \\
$y$ of $v_{\max}$ ($0.447\, Y_{\Ro}$) & $1.14\times 10^6$~m & $1.18\times 10^6$~m & 3.5\% \\
cell edge $Y_{\Ro}$ (vs jet latitude) & $2.55\times 10^6$~m & $2.36\times 10^6$~m & $-7.4\%$ \\
\bottomrule
\end{tabular}
\end{center}

Sections \ref{sec:edge}--\ref{sec:v} upgraded to \textbf{[confirmed at the
1--8\% level, single run]} for these observables at this $\Ro$. Additional
findings: the $u$ plateau outside the cell is the parabolic forcing's RCE
gradient wind,
\begin{equation}
  u_{\mathrm{RCE}} = \frac{2 g H \Delta_T}{T_0 \beta y_1^2} = 18.35~\text{m/s},
\end{equation}
constant out to $y_1$ (verified to 3 digits); the $v$-threshold edge
diagnostic is unusable at small $v_d$ because the drag circulation
$v = \epsilon_u u_{\mathrm{RCE}}/(\beta y) \sim 2\times 10^{-3}$~m/s
exceeds 10\% of $v_{\max}$ out to $\sim 9.3\times 10^6$~m; the
$\theta$-merge diagnostic must search poleward of the jet, not of the $v$
max, to avoid the equal-area lobe crossing.

\subsection*{2026-07-16, check 1b: $v_d = 0$, vertical advection off}

Config: as 1a plus \texttt{-{}-no-vert-advec-u}; converged day 971.
Output: \path|.../tier0_amc_edge/amc_ss09_ny801_dt30_vd0_novert.nc|.

Removing vertical advection left the cell-mean $\Ro$ essentially unchanged
(fit 0.473 vs 0.470; area mean 0.536 vs 0.469), because \textbf{the
Rayleigh drag, at its default $\epsilon_u = 10^{-8}$~s$^{-1}$, is the
leading-order non-AMC agent for this forcing} \textbf{[computed]}: the
steady zonal momentum balance over the cell interior is
\begin{equation}
  v \left( \beta y - \dpar{u}{y} \right) = \epsilon_u u
  \quad \text{to within ${\sim}5$--$10\%$}
\end{equation}
(median ratio 1.05 over $0.3 \times 10^6 < y < 2.2\times 10^6$~m). The
parabolic forcing's weak circulation ($v_{\max} \sim 9\times
10^{-3}$~m/s) makes $(1 - \Ro) \sim \epsilon_u u / (v \beta y)$ order
one. Local $\Ro(y)$ declines monotonically from ${\sim}1$ at the ascent to
${\sim}0.2$ at the jet.

Theory at diagnosed $\Ro = 0.473$ vs run \textbf{[computed]}: depression
0.378 predicted vs 0.358~K measured (5.3\%); $v_{\max}$ $9.18\times
10^{-3}$ vs $8.81\times 10^{-3}$~m/s (4.2\%); edge $2.54\times 10^6$ vs
jet $2.36\times 10^6$~m ($-7.1\%$).

Consequences for the suite: (i) the $\Ro \to 1$ anchor needs reduced
$\epsilon_u$, so the Tier-3 drag ladder is load-bearing rather than
optional; (ii) at the sin$^2$ forcing the circulation is ${\sim}3\times$
stronger, so the drag share is ${\sim}3\times$ smaller but still
non-negligible at $v_d = 0$ (consistent with Zhang et al.\ 2025 Fig.~4's
max $\Ro \sim 0.9$ there); (iii) the steady-state detector (window 30~d,
$10^{-4}$) fires while the jet still drifts ${\sim}0.4\%$ per 60~d on the
drag timescale, so budget-grade runs need a stricter threshold or a fixed
post-detection extension.

\subsection*{2026-07-16, profile-level addendum to runs 1a--1b}

The tables above are endpoint checks (the equatorial depression and
$v_{\max}$ are, respectively, the $y = 0$ anchor of $T$ and a single point
of $v$). The suite standard is now the full-profile scorecard
(Section~\ref{sec:diag}), applied retroactively to both runs: final-30-day means, theory at the
fitted cell-mean $\Ro$ of 0.479 (1a) and 0.473 (1b). Under this
standardized window 1b's fit reproduces the value quoted above exactly,
while 1a's shifts from 0.470 to 0.479, within run 1a's stated ${\sim}2\%$
drift allowance. Profile results \textbf{[computed]}:

\begin{center}
\begin{tabular}{llcccc}
\toprule
run & field & $\max|\Delta|$ (cell) & at $y$ [$10^6$ m] & \% of scale
  & inner half \\
\midrule
1a & $u$ & 10.4~m/s & 2.52 & 34 & 2.3~m/s (7.5\%) \\
1a & $T$ & 0.133~K & 2.01 & 36 & 0.038~K (10\%) \\
1a & $v$ & 4.0~mm/s & 2.44 & 45 & 0.49~mm/s (5.5\%) \\
1b & $u$ & 10.2~m/s & 2.52 & 33 & 3.2~m/s (10\%) \\
1b & $T$ & 0.152~K & 2.01 & 40 & 0.042~K (11\%) \\
1b & $v$ & 3.9~mm/s & 2.44 & 43 & 0.48~mm/s (5.2\%) \\
\bottomrule
\end{tabular}
\end{center}

Scales: cell-edge $u_{\Ro}(Y_{\Ro})$, the equatorial depression, and
$v_{\max}$. Reading: the inner half of the cell agrees at the 5--11\%
level in all three fields, and the misfit concentrates at the terminus in
all three, consistent with the diagnosed local $\Ro(y)$ declining to
${\sim}0.2$ at the jet: a single cell-mean $\Ro$ is furthest from the
local value exactly there. The simulated subtropical warm lobe of $D(y)$
peaks at 0.17~K (1a) against the predicted 0.30~K and is smeared poleward
across the edge ($D$ stays ${\sim}{+}0.015$~K out to at least $4.5\times
10^6$~m), suggesting the sharp-edge equal-area closure, not the quartic
interior shape, is the limiting idealization. The $u$ misfit sits at the
theory's terminus discontinuity (the jump to the RCE plateau). The ms's
linear-$\Ro$ variant targets exactly this regime and is testable against
these same stored profiles (suite extension, not yet scheduled).

\end{document}
